\section{Постановка задачи}
\label{sec:Chapter1} \index{Chapter1}

\section{Цели работы:}

\begin{enumerate}
    \item Расширить типовую систему верификатора для более точного отображения типовой
    системы виртуальной машины на типовую систему верификатора.
    \item Разработать и внедрить алгоритмы проверок ряда инструкций для расширения множества
    верифицируемых программ.
\end{enumerate}

\section{Задачи работы:}
\begin{enumerate}
    \item Изучить существующие решения.
    \item Сравнить типовые системы виртуальной машины и верификатора, выявить неверифицируемые типы.
    \item Расширить типовую систему и поддержать логику обработки данных типов.
    \item Проанализировать архитектуру набора команд платформы и выявить множество неверифицируемых
    инструкций.
    \item Реализовать и внедрить алгоритм верификации новых инструкций.
    \item Протестировать каждый из предложенных алгоритмов.
    \item Внедрить предложенное решение в программный продукт.

\end{enumerate}

\subsection{Ожидаемый результат}

Ожидаемым результатом работы является расширение множества верифицируемых программ, которые могут быть
проверены до исполнения, а также повышение точности анализа в тех
случаях, где прежняя типовая система была недостаточно выразительной. Практически это должно
привести к уменьшению числа ложных отказов верификации, более раннему обнаружению ошибок
кодогенерации и повышению надёжности механизмов исполнения для ETS-специфичных и инструкций, связанных
с взаимодействием с динамическими средами.

\subsection{Критерии достижения цели}

Цель работы можно считать достигнутой, если верификатор получит поддержку новых инструкций,
будет корректно обрабатывать ранее проблемные типовые случаи и подтверждать это на наборе добавленных тестов. 
Дополнительным критерием является практическая значимость внесённых
изменений: способность выявлять ошибки компилятора и улучшать гарантии платформы по безопасности и
корректности исполнения. И если по итогам работы верификация станет включенной по умолчанию в основной открытой ветке
разработки платформы, то это будет свидетельствовать о полезности и успешности.
